\documentclass[11pt, landscape]{article}

% change the size of the page
\usepackage[
	papersize={20in, 11.5in},
	margin=1in
]{geometry}

\usepackage[table, dvipsnames]{xcolor}
\usepackage{tabularx}
\usepackage{fontspec}
\usepackage{etoolbox}
\usepackage{enumitem}
\usepackage{environ}
\usepackage{xltabular}
\usepackage[most]{tcolorbox}

\setmainfont{PT Sans} 
\usepackage[portuguese]{babel}

% --- Color Palette ---
\definecolor{primaryBlue}{HTML}{2C3E50}
\definecolor{normalGreen}{HTML}{27AE60}
\definecolor{altOrange}{HTML}{E67E22}
\definecolor{excRed}{HTML}{C0392B}
\definecolor{lightGray}{HTML}{F8F9FA}
\definecolor{borderGray}{HTML}{D1D1D1}

% --- Table Styling ---
\renewcommand{\arraystretch}{1.6}
\setlength{\arrayrulewidth}{0.5pt}
\arrayrulecolor{borderGray}

% --- Counters ---
\newcounter{stepcount}
\newcounter{substepcount}

% --- Commands ---

\newcommand{\ucHeader}[5]{
	\noindent {\color{primaryBlue}\Huge \textbf{Caso de Uso: #1}}
	
	\vspace{1.5em}
	\noindent
	\begin{tabularx}{\textwidth}{|l|X|}
		\hline
		\rowcolor{lightGray} \textbf{Descrição} & #2 \\ \hline
		\textbf{Cenários} & #3 \\ \hline
		\rowcolor{lightGray} \textbf{Pré-condições} & #4 \\ \hline
		\textbf{Pós-condições} & #5 \\ \hline
	\end{tabularx}
}

\NewEnviron{flowtable}{
	\noindent
	\begin{xltabular}{\textwidth}{|>{\hsize=1.2\hsize}X|>{\hsize=0.8\hsize}X|>{\hsize=1.4\hsize}X|>{\hsize=0.6\hsize}X|}
		\hline
		\rowcolor{lightGray} 
		\textbf{Passos} & \textbf{Lógica de Negócio} & \textbf{Definição da API} & \textbf{Subsistema} \\ \hline
		\endfirsthead
		
		\hline
		\rowcolor{lightGray} 
		\textbf{Passos} & \textbf{Lógica de Negócio} & \textbf{Definição da API} & \textbf{Subsistema} \\ \hline
		\endhead
		
		\BODY
	\end{xltabular}
}

\newcommand{\step}[4]{
	\stepcounter{stepcount}
	\thestepcount. #1 & #2 & \texttt{#3} & #4 \\ \hline
}

\newcommand{\substep}[4]{
	\stepcounter{substepcount}
	\theparentstep.\thesubstepcount\ #1 & #2 & \texttt{#3} & #4 \\ \hline
}

% lazy step: use it when the operation will not be reflected in the system
\newcommand{\lstep}[1]{
	\step{#1}{-}{-}{-}
}

% lazy sub step: use it when the operation will not be reflected in the system
\newcommand{\lsubstep}[1]{
	\substep{#1}{-}{-}{-}
}

\newcommand{\normalFlowHeader}{
	\setcounter{stepcount}{0}
	
	\vspace{1em}
	
	\noindent
	\begin{tcolorbox}[
		enhanced,
		sharp corners,
		boxrule=0pt,
		frame hidden,
		colback=normalGreen!8!white,
		left=5pt, right=5pt,
		top=8pt, bottom=8pt,
		after skip=0pt,
		]
		{\color{normalGreen}\Large\bfseries Fluxo Normal}
	\end{tcolorbox}
	\nointerlineskip
}

\newcommand{\subFlowHeader}[4]{
	\setcounter{substepcount}{0}
	\def\theparentstep{#3}
	
	\noindent
	\begin{tcolorbox}[
		enhanced,
		sharp corners,
		boxrule=0pt,
		frame hidden,
		colback=#1!8!white,
		left=5pt, right=5pt,
		top=5pt, bottom=5pt,
		after skip=0pt,
		]
		{\color{#1}\Large\bfseries #2} \vspace{0.3em} \\
		\noindent\hspace{2pt}\vrule width 2pt depth 0.8ex height 2.2ex \hspace{0.5em}
		\begin{minipage}{\dimexpr\textwidth-25pt}
			\small \textbf{PASSO #3} \quad \textbullet \quad \textbf{CONDIÇÃO:} \textit{#4}
		\end{minipage}
	\end{tcolorbox}
	\nointerlineskip
}

\begin{document}

	\ucHeader{Pagar Multas}
	{Funcionário regista o pagamento de multas pelo atraso na entrega de livros requisitados}
	{Cenários 2 e 5}
	{Funcionário autenticado}
	{Multas Pagas}
	
	\normalFlowHeader
	\begin{flowtable}
		\step{Sistema calcula valor de multas a pagar}{Calcular o valor de multas a pagar para o utente}{calcularValorMultas(id : String) : float}{Utentes}
		\lstep{Funcionário confirma pagamento de multas}
		\step{Sistema imprime comprovativo de pagamento de multas}{Gerar comprovativo de pagamento de multas}{gerarComprovativoPagamento(codPag : String) : Comprovativo}{Utentes}
	\end{flowtable}
	
	\subFlowHeader{excRed}{Fluxo de Exceção}{2}{Utente não paga as multas}
	\begin{flowtable}
		\lsubstep{Sistema regista que utente não pagou as multas}
	\end{flowtable}
	
	\newpage
	
	\ucHeader{Validar Credenciais de Utente}
	{Sistema valida credenciais de utente para efetuar requisições, entregas e pagamento de multas}
	{Cenário 1, 2, 3, 4 e 5}
	{True}
	{Credenciais autenticadas}
	
	\normalFlowHeader
	\begin{flowtable}
		\lstep{Funcionário apresenta número e nome de utente}
		\step{Sistema valida número e nome de utente}{Validar credenciais de utente}{validarUtente(id : String, nome : String) : boolean}{Utentes}
	\end{flowtable}
	
	\subFlowHeader{excRed}{Fluxo de Exceção}{2}{Número e nome inexistentes}
	\begin{flowtable}
		\lsubstep{Sistema indica que utente é inexistente}
	\end{flowtable}
	
	\newpage
	
	\ucHeader{Registar Entrega de Livro}
	{Funcionário regista a entrega de livros requisitados por utentes}
	{Cenário 4 e 5}
	{Funcionário autenticado}
	{Devolução do livro registada e estado do livro alterado para 'Disponível'}
	
	\normalFlowHeader
	\begin{flowtable}
		\lstep{\texttt{<<include>>} Validar Credenciais de Utente}
		\lstep{Funcionário fornece identificação do livro}
		\step{Sistema valida identificação do livro}{Verificar se livro pertence à biblioteca}{livroExiste(id : String) : boolean}{Livros}
		\step{Sistema valida que entrega está dentro do prazo}{Verificar se entrega está dentro do prazo}{entregaDentroPrazo(utente : String) : boolean}{Utentes}
		\step{Sistema regista a entrega do livro}{Registar entrega de livro}{registarEntrega(idLivro : String, idUtente : String) : Devolucao}{Livros}
		\step{Sistema altera estado do livro para 'Disponível'}{Alterar estado de livro}{alterarEstado(livro : String, estado : String) : void}{Livros}
		\step{Sistema gera comprovativo de devolução}{Gerar comprovativo de devolução}{gerarComprovativoDevolucao(idDevol : String) : Comprovativo}{Livros}
	\end{flowtable}
	
	\subFlowHeader{altOrange}{Fluxo Alternativo}{4}{Entrega está fora do prazo}
	\begin{flowtable}
		\lsubstep{\texttt{<<include>>} Pagar Multas}
		\lsubstep{Regressa a 5}
	\end{flowtable}
	
	\subFlowHeader{excRed}{Fluxo de Exceção}{3}{Livro não pertence à biblioteca}
	\begin{flowtable}
		\lsubstep{Sistema informa que livro não pertence à biblioteca}
	\end{flowtable}
	
	\newpage
	
	\ucHeader{Registar Requisição de Livro}
	{Funcionário regista a requisição de um livro por parte de um utente}
	{Cenário 1, 2 e 3}
	{Funcionário autenticado}
	{Registo de requisição adicionado ao sistema e estado do livro alterado para 'Requisitado'}
	
	\normalFlowHeader
	\begin{flowtable}
		\lstep{\texttt{<<include>>} Validar Credenciais de Utente}
		\step{Sistema valida que utente não tem multas a pagar}{Verificar se utente não tem multas a pagar}{multasPorPagar(utente : String) : boolean}{Utentes}
		\lstep{Funcionário indica código do livro}
		\step{Sistema verifica disponibilidade do livro para ser requisitado}{Verificar a disponibilidade do para ser requisitado}{podeSerRequisitado(livro : String) : boolean}{Livros}
		\step{Sistema regista requisição do livro}{Registar requisição de livro}{registarRequisicao(livro : String, utente : String) : Requisicao}{Livros}
		\step{Sistema altera estado do livro para 'Requisitado'}{Alterar estado do livro}{alterarEstadoLivro(livro : String, estado : String) : void}{Livros}
		\step{Sistema calcula data de devolução do livro}{Calcular data de devolução de livro}{getDataDevolucao(livro : String) : LocalDateTime}{Livros}
		\step{Sistema gera comprovativo da requisição}{Gerar comprovativo de requisição}{gerarComprovativoRequisicao(requisicao : String) : ComprovativoRequisicao}{Livros}
	\end{flowtable}
	
	\subFlowHeader{altOrange}{Fluxo Alternativo}{2}{Utente tem multas a pagar}
	\begin{flowtable}
		\lsubstep{\texttt{<<include>>} Pagar Multas}
		\substep{Sistema prolonga a data de devolução dos livros por entregar}{Prolongar data de devolução de livros}{prolongarDataDevolucao(utente : String) : void}{Utentes}
		\lsubstep{Regressa a 3}
	\end{flowtable}
	
	\subFlowHeader{excRed}{Fluxo de Exceção}{4}{Livro não pode ser requsitado}
	\begin{flowtable}
		\lsubstep{Sistema informa que livro não pode ser requisitado}
	\end{flowtable}
	
\end{document}